\section{引言}

\subsection{国家战略与发展需求}

建设数字中国是数字时代推进中国式现代化的重要引擎，也是构筑国家竞争新优势的重要支撑。《数字中国建设整体布局规划》明确提出，要夯实数字基础设施和数据资源体系，强化数字技术创新能力，全面提升数字中国建设的整体性、系统性和协同性\footnote{政策来源：中共中央、国务院，《数字中国建设整体布局规划》，2023年2月；\href{https://www.gov.cn/zhengce/2023-02/27/content_5743484.htm}{中国政府网原文}。}。《关于深化智慧城市发展 推进城市全域数字化转型的指导意见》进一步将城市作为推进数字中国建设的综合载体，要求整合状态感知、建模分析、城市运行和应急指挥等能力，实现城市态势全面感知、趋势智能研判、协同高效处置和调度敏捷响应\footnote{政策来源：国家发展改革委、国家数据局、财政部、自然资源部，《关于深化智慧城市发展 推进城市全域数字化转型的指导意见》（发改数据〔2024〕660号），2024年5月；\href{https://www.ndrc.gov.cn/xxgk/zcfb/tz/202405/t20240520_1386326_ext.html}{国家发展改革委原文}。}。国务院《关于深入实施“人工智能+”行动的意见》则面向科技、产业、民生和治理等重点领域，提出推动人工智能与经济社会各行业各领域广泛深度融合，发展新一代智能终端和智能体，提升公共安全、防灾减灾、交通治理和生产运行等场景的智能化水平\footnote{政策来源：国务院，《关于深入实施“人工智能+”行动的意见》（国发〔2025〕11号），2025年8月；\href{https://www.cac.gov.cn/2025-08/27/c_1758018277755538.htm}{中央网信办转载的中国政府网原文}。}。上述国家战略部署表明，城市数字化转型正在从基础设施联网和数据汇聚，进一步迈向全域感知、智能理解与协同决策的新阶段。

与此同时，新型城市基础设施建设、韧性城市建设和新型工业化对感知技术提出了更具体的任务要求。国家层面的城市数字化转型部署已明确提出，推进城市公共设施数字化改造和智能化运营，统筹部署泛在、韧性的城市智能感知终端，并推动城市智能基础设施与智能网联汽车协同发展。城市道路、建筑、市政设施、工业园区和公共空间由此加快部署物联网终端、智能摄像头、雷达、无人机、移动机器人及车路协同设施，城市运行逐步由周期性人工巡查转向连续在线监测，由事后处置转向风险预测与主动干预。低空飞行、智能网联汽车、智能制造等新产业与新场景进一步拓展了感知对象的速度范围、尺度范围和空间范围：系统既要识别快速运动的车辆和飞行器，也要发现远距离细线、微小缺陷和弱小目标；既要在白天和常规天气下稳定运行，也要在暗光、雨雪、遮挡和强动态干扰下持续工作；既要支撑单一设备的实时控制，也要服务跨区域、多节点、多主体的协同行动\footnote{政策依据：发改数据〔2024〕660号第（八）项。}。

这些部署共同指向一个基础性问题：数字城市首先必须能够真实、连续、完整地感知物理城市。如果前端观测在极端条件下失真，后续算法就无法恢复已经丢失的关键状态；如果多源信息不能在有限通信资源下有效流动，再多的感知节点也难以形成全局认知；如果复杂任务只能依赖高成本模型进行集中式端到端计算，系统也难以满足城市运行所要求的低时延和高可靠性。因此，感知技术已不再只是摄像头或传感器的数据采集问题，而是贯穿物理信号获取、跨源信息组织、任务语义理解和行动决策的系统性问题。提升这一链条的极限感知能力、全域协同能力和复杂认知能力，是数字中国、智慧城市与“人工智能+”战略由基础设施建设走向实际治理效能的重要技术基础。

从数字化转型的演进过程看，城市智能化大体经历“设施数字化、数据资源化、业务智能化、决策协同化”的能力递进。设施数字化解决物理对象能否接入网络的问题，数据资源化解决不同部门、不同主体的数据能否汇聚和流通的问题，业务智能化解决数据能否支撑识别、预测与处置的问题，决策协同化则要求多个平台、多个智能体围绕共同目标形成实时联动。感知处于这条链路的起点，也决定着后续环节的能力上限：只有前端信息真实、及时且具有足够覆盖度，数据才能被有效组织，模型推理才能建立在可靠事实上，最终决策才具有可解释、可执行和可追溯的基础。

这一基础作用在典型城市应用中表现得尤为突出。在智慧交通与低空运行场景中，感知结果直接关联目标跟踪、航迹预测、冲突检测和运动控制，系统不仅要“识别到”，还要以足够高的更新频率持续输出状态；在智能制造与基础设施运维场景中，微小裂纹、细线结构和早期故障征兆往往尺度小、对比度低，需要在复杂背景下保留弱而关键的异常信息；在智能安防与应急处置场景中，暗光、烟雾、雨雪、遮挡和通信波动可能同时出现，系统还必须整合固定摄像头、移动终端、无人机与机器人等多方观测，形成对事件演化过程的整体判断。不同场景虽然任务对象各异，但共同要求感知体系由“单设备获得局部图像”升级为“多节点持续获得可信状态，并服务后续认知与行动”。

\subsection{城市复杂场景的感知需求与关键挑战}

智能感知是连接物理城市与数字城市的基础环节，也是支撑智慧交通、智能制造、智能安防和智慧应急等产业应用的共性关键技术。只有持续、准确地获取城市运行状态，并将分散、多源、异构的观测转化为可理解、可推理、可行动的信息，城市智能系统才能形成可靠的分析、决策和控制闭环。随着城市应用不断向高速动态、极端环境、广域覆盖和复杂任务延伸，感知技术所面对的需求呈现出三个显著趋势：一是需求极端化，感知对象由常规目标扩展至高速、微小和暗光目标；二是规模扩大化，感知范围由单点、单域扩展至空天地海多域以及海量节点协同；三是任务复杂化，系统能力由目标检测和几何测量扩展至复杂语义理解、空间认知与系统规划。传统感知体系在这些新需求下逐渐暴露出“看不清、看不全、看不透”的瓶颈。

上述需求变化并不是对传统感知指标的简单提高，而是带来了传感机理、信息组织方式和计算范式的共同变化。从项目所面向的典型应用看，关键挑战可以归纳为以下三个相互关联的层次。

\subsubsection{极端条件下单点感知能力不足}

在城市交通和低空飞行中，目标运动速度快、状态变化频繁，传统RGB帧相机受固定曝光和有限帧率制约，容易产生运动模糊并遗漏帧间变化，感知输出频率往往难以匹配飞行控制、避障和应急响应所需的高频状态反馈。在工业质检、基础设施巡检和安防监控中，待检测对象可能只占据少量像素，甚至小于单个像素的等效尺度；其边缘弱、纹理少，并容易与复杂背景混叠，传统基于外观纹理的识别方法难以稳定提取判别特征。在夜间安防、地下空间、雨雪天气和强明暗交替场景中，传统相机的动态范围和信噪比进一步下降，有效目标信号可能被背景噪声和环境干扰淹没。

这类问题表面上分别表现为高速目标难跟踪、微小目标难识别和暗光目标难感知，本质上则来自传统帧式成像对时间、空间和光照信息的统一采样方式。提高帧率通常会缩短曝光时间并降低单帧信噪比，提高空间分辨率又会增加数据量和计算负担，延长曝光则会加剧高速运动模糊，单纯通过堆叠像素、提高帧率或增大模型规模难以从根本上消除这些制约。因此，需要从感知机理出发，引入对亮度变化和运动信息进行异步响应的新型视觉传感方式，并针对事件信号的稀疏、异步和高时间分辨特性重新设计检测、定位与增强算法。

\subsubsection{大规模场景下全域协同观测不完备}

城市感知正在由单个摄像头、单台车辆和单架无人机，扩展到由固定设施、移动终端、边缘节点和云端平台共同组成的分布式网络。单个节点受视场、遮挡、量程和部署位置限制，只能获得局部观测；视觉、事件、雷达、惯性和点云等不同模态虽然能够提供互补信息，却具有不同的采样频率、空间坐标、数据密度与噪声规律。若缺少统一的时空表征与匹配机制，多种传感器的简单叠加不仅不能形成互补，反而可能因时间错位、空间偏差和特征冲突降低融合精度。

当节点数量和数据规模进一步增加时，信息交换本身也成为系统瓶颈。原始图像、点云和神经网络中间特征具有较高数据维度，在带宽有限、链路动态和端侧算力受限的条件下持续传输，会引发通信拥塞、处理延迟和能源消耗。与此同时，不同观测的价值和可靠性并不相同：部分节点掌握关键视角和关键时刻的信息，部分数据则高度重复；部分传感器在当前环境中可靠，另一些传感器可能因遮挡、漂移或恶劣天气产生失真。如果系统不能回答“不同信息如何互补、哪些信息值得传输、哪些观测应当采信”，就难以将多源数据真正转化为完备、实时、可信的全域观测。

\subsubsection{复杂任务中认知准确性与实时性难以兼顾}

感知系统的任务正在由目标类别、位置和轨迹等几何层信息，向事件关系、场景因果、空间语义、任务分解和协同规划等高层认知延伸。例如，城市应急系统不仅要发现烟雾、车辆和人群，还要区分正常生产与异常火情，结合道路拥堵状态规划救援路径；多无人机和无人车协同搜索不仅要识别目标，还要理解自然语言任务、维护全局空间记忆，并根据平台能力和实时观测持续调整分工。这些任务包含多步推理和长时程决策，依靠单一轻量模型往往难以获得足够的语义理解与泛化能力。

基础大模型和视觉语言模型为复杂认知提供了新的能力，但其高质量推理通常伴随较高的词元、显存、算力和时间开销。现实任务中，大量工作只是高频、局部、规则明确的感知与执行操作，并不需要每次调用大模型；若将简单任务、复杂任务、全局规划和局部控制全部集中到大模型中处理，会造成计算资源错配，并使推理延迟阻塞实时执行。另一方面，完全依靠小模型又难以处理开放语义、长链条任务和陌生环境。因此，需要明确大小模型、全局与局部规划模块以及不同智能体之间的能力边界，依据任务难度、空间尺度和信息价值进行动态认知卸载，使高阶模型集中处理真正需要复杂推理的问题。

\subsection{传统技术路线的局限与本项目的研究切入点}

面向上述需求，现有研究通常沿着提高传感器规格、扩大模型规模或增加感知节点数量三条路线提升系统性能，这些路线在常规场景中具有明确作用，但难以单独应对城市复杂环境中的多重约束。首先，提高图像分辨率和采样帧率会同步增加数据吞吐、存储和计算负担，而且固定曝光成像仍然受到运动模糊、光照变化和帧间信息缺失的制约。换言之，单纯增加数据量并不必然增加有效信息量；当真正决定任务的高速变化、弱小边缘或瞬态信号在采集阶段已经丢失时，后端模型规模的增长无法完整补回物理观测中不存在的信息。

其次，增加传感器和节点数量能够扩展覆盖范围，却也把问题由“信息不足”转化为“信息过量且质量不一”。全量上传原始数据、在云端进行集中融合的方式依赖稳定而充足的通信条件，难以适应移动平台、突发事件和大规模并发接入；对所有节点采用固定权重或静态融合规则，又无法反映遮挡、漂移、恶劣天气和任务变化造成的可靠性波动。全域感知因而不能仅以节点数量衡量，还应同时考察异构信息是否对齐、通信资源是否用在关键观测上，以及融合结果能否识别和抑制不可信信息。

再次，将复杂感知、语义理解和任务规划统一交给大型模型，虽然能够减少人工设计的中间环节，却可能把高频简单操作也纳入高成本推理过程。城市智能系统通常同时包含毫秒至秒级的状态更新、分钟级的任务调整以及更长时间尺度的全局规划，不同环节对于语义能力、计算精度和响应速度的要求并不相同。若缺少分层机制，高层推理可能阻塞底层控制，底层大量重复信息也会挤占高层模型的上下文与算力；若完全拆分，各模块之间又可能因缺少共享语义和任务反馈而形成新的信息孤岛。

基于此，本项目不把感知、通信、融合和推理视为相互独立的算法模块，而是以有效信息在系统中的产生、流动与利用为主线，从三个层面形成研究切入点：在信息产生阶段，利用动态视觉的变化触发机制提高极端场景中关键物理信号的获取效率；在信息流动阶段，利用多感官互补、选择传导和先验调节机制建立多源协同观测；在信息利用阶段，利用认知卸载机制把任务合理分配至大小模型、不同空间层级和异构智能体。三者共同解决“有限传感资源如何获得更关键的信息、有限通信资源如何传递更有价值的信息、有限计算资源如何完成更复杂的任务”这一贯穿全链路的核心矛盾。

\subsection{生物启发式感知的研究依据}

生物感知系统为突破上述瓶颈提供了重要启发。生物视觉通过不同感光与神经通路分工，实现对细节、运动和弱光信息的高效编码；生物多感官系统通过感官互补、选择传导和经验调节，在有限神经资源下形成完整而可信的环境认知；生物认知系统则通过快慢通路、记忆分层以及大脑与感觉运动系统之间的协作，将高频简单任务与低频复杂任务分开处理。借鉴这些机理，可以从传感、融合到推理重新组织机器感知过程，使有限的感知、通信和计算资源聚焦于最具任务价值的信息，从而拓展城市智能系统的感知极限与认知边界。

在视觉感知层面，生物视觉由不同类型的感光细胞和神经通路分别处理颜色、细节、亮度与运动变化。事件相机借鉴生物视网膜对变化刺激的响应方式：每个像素独立响应亮度变化，只在变化超过阈值时输出事件，从而避免重复记录静态背景，并以微秒级时间分辨率保留高速运动过程。由此，机器视觉可以由“按固定帧率完整记录场景”转向“按变化选择性记录信息”，为突破高速、微小和暗光场景中的成像限制提供新的传感基础\footnote{科学依据：G. Gallego等，``Event-based Vision: A Survey,'' IEEE TPAMI, 44(1):154--180, 2022，DOI: \href{https://doi.org/10.1109/TPAMI.2020.3008413}{10.1109/TPAMI.2020.3008413}。}。

在多源融合层面，生物体不会依赖单一感官形成对环境的全部判断。视觉能够描述空间结构和外观，听觉能够在视线受阻时提供方向和时序线索，触觉与本体感觉则能够反馈接触状态和自身运动。更重要的是，神经系统会根据任务、刺激强度和既有经验动态分配注意力：高价值刺激被优先传导，相互重复或无关的信息受到抑制，不同感官的可信程度也会随环境条件调整。这一机制提示，多传感器、多节点系统的关键并非无差别汇聚数据，而是首先形成跨模态互补，进而根据任务价值选择信息，最后结合先验知识判断信息是否可信。感官互补、选择传导与先验调节共同构成了全域协同融合的仿生依据。

在认知推理层面，人类可以通过外部动作、环境表征和工具改变任务的信息处理需求，降低内部认知负担，这一现象通常被概括为认知卸载。对应到机器系统，本项目将这一原则进一步抽象为模型与智能体之间的任务分工：将高频简单任务分配给轻量模型和局部智能体，将低频复杂任务交由基础大模型和全局规划模块，并利用分层记忆、空间分解和视觉词元筛选进一步减少冗余推理，使有限算力与任务复杂度相匹配\footnote{概念来源：E. F. Risko与S. J. Gilbert，``Cognitive Offloading,'' Trends in Cognitive Sciences, 20(9):676--688, 2016，DOI: \href{https://doi.org/10.1016/j.tics.2016.07.002}{10.1016/j.tics.2016.07.002}。}。

本项目所强调的“生物启发”并非对生物器官结构的机械复刻，而是提炼其中具有普适意义的信息处理原则：以变化触发代替均匀采样，以功能互补代替单源增强，以价值选择代替全量传输，以先验调节代替静态加权，以分层卸载代替集中计算。上述原则分别作用于感知信息的产生、传输、融合和理解过程，能够将三个看似不同的技术问题组织在统一的方法论之下。

\subsection{研究对象与总体思路}

本项目面向城市复杂场景，研究从物理信息获取、多源信息协同到复杂任务理解的生物启发式感知理论、方法与系统。研究对象不是单一传感器或单一识别算法，而是贯穿“感知接入---协同融合---认知推理”的完整信息处理链条：前端关注高速、微小和暗光目标能否被及时、清晰地观测；中端关注跨模态、跨时空和跨节点信息能否高效、可信地组织与融合；后端关注复杂语义、空间关系和系统任务能否在有限算力下准确、实时地理解与规划。

围绕上述研究对象，本项目以“让机器看得更清、看得更广、看得更透”为总体目标，以生物动态视觉、多感官协同和认知卸载机理为共同方法论，形成“单点感知能力增强---多点协同感知体系增强---复杂任务理解能力增强”的总体技术路线。其中，“看得更清”旨在突破传统成像在时间分辨率、空间分辨率和动态范围上的限制；“看得更广”旨在突破单一模态、单一节点和静态采信机制对全域观测的限制；“看得更透”旨在突破复杂任务对大模型、海量词元和集中式计算的过度依赖。三类研究相互衔接，分别回答“信息如何被有效获取”“信息如何被高效而可信地协同”“信息如何被准确而实时地理解”三个核心问题。

从基础研究角度看，本项目重点研究三类共性问题。第一，研究非均匀、异步脉冲信息在高速动态和低信噪比条件下的有效表征规律，探索时间分辨率、空间细节和环境鲁棒性之间新的协调方式；第二，研究异构模态和分布式节点之间的信息对应关系、任务价值度量与可信度演化规律，使多源观测能够在通信和计算约束下形成稳定一致的场景表达；第三，研究复杂任务在模型能力、空间层级和执行主体之间的合理分解机制，使语义理解、记忆规划和实时控制能够以不同时间尺度协同运行。这三类问题从信号表征、信息融合延伸到认知计算，构成本项目的主要研究边界。

从系统研究角度看，本项目强调算法、传感器、通信网络、边缘计算与智能体执行平台之间的协同设计。前端传感算法需要保留事件信号的高频与稀疏优势，中端融合算法需要适应真实网络中的带宽波动和节点异质性，后端推理算法需要考虑端侧设备的算力、时延和任务安全要求；各层输出既是下一层的输入，也会受到最终任务反馈的调节。通过这种端到端但非单体化的系统设计，项目力求避免局部算法指标提升与整体应用效果脱节，使研究成果能够部署到无人机、车辆、机器人、固定感知设施和城市级平台之中。

\subsection{总体研究目标与技术路线}

本项目的总体研究目标，是构建面向城市复杂场景的生物启发式智能感知技术体系，在传感条件受限、观测来源异构、通信资源有限和认知任务复杂的情况下，仍能够稳定完成环境观测、信息融合与任务推理。该目标并非追求某一个孤立指标的提升，而是强调信息完整性、处理实时性、系统可靠性与资源可承受性之间的综合协调。针对不同技术层级，项目分别设置极限场景有效感知、分布式信息高效协同和复杂任务分层推理三类目标，并通过统一的信息价值导向把三类目标连接起来。

在前端感知层，技术路线由传统帧式图像的均匀采样扩展到事件信号的异步变化采样，并结合帧图像、惯性或空间结构等互补观测，围绕高速、微小、暗光和恶劣天气条件设计时空表征方法。其关键不是简单替代传统相机，而是充分利用不同传感机制各自保留的信息：以事件流表达快速变化和精细运动边缘，以帧或结构信息提供稳定外观和全局参照，再依据目标尺度、运动状态和环境条件组织多时间尺度信息。前端输出由单一图像识别结果扩展为连续、可定位、可传递的目标与环境状态，为后续协同融合提供高质量输入。

在中端协同层，技术路线由“全量数据集中上传”转向“跨模态对齐、任务价值选择和可信先验调节”。系统首先建立事件、视觉、惯性和三维结构等异构信息之间的时空对应关系，再根据节点观测质量、任务相关性和信息新颖度确定传输与融合优先级，最后结合历史经验、语义关系、空间结构和运动规律校正失真观测。通过这一过程，多节点系统既能够利用不同视角和不同模态弥补单点盲区，又能够降低重复信息对通信和计算资源的占用，并在个别节点失效或环境变化时保持融合结果的稳定性。

在后端认知层，技术路线由单一模型集中推理转向“任务难度分流、空间层级分解和规划执行协同”。轻量模型承担频繁出现、边界明确且对时延敏感的局部感知与控制任务，基础大模型处理开放语义理解、长链条推理和全局任务规划；全局记忆维护长期目标与整体空间关系，局部模块面向当前可见区域完成即时行动；高阶模型给出任务分解和协作策略，车辆、无人机、机器人等执行主体依据自身能力并行响应。通过分层调用和闭环反馈，系统在保留大模型复杂认知能力的同时，减少不必要的高成本推理，使认知精度与执行实时性能够共同进入设计过程。

项目的技术评价也与上述跨层路线相对应。前端不仅关注检测与定位精度，还关注高速响应、弱信号保留和环境适应能力；中端除融合效果外，还考察通信开销、信息时效、节点扩展性和异常观测下的稳健性；后端则综合考察任务完成质量、推理时延、计算资源消耗和多智能体协同效率。最终通过典型城市应用验证各层方法之间的接口与反馈机制，形成从关键算法、功能模块到系统平台的递进式成果，使理论方法能够在真实设备、真实网络与真实任务约束下接受检验。

\subsection{主要研究内容}

\textbf{（1）动态视觉启发的极限场景感知。}
针对城市交通、低空飞行、工业检测、智能安防和应急处置等场景中高速目标难跟踪、微小细粒度目标难识别、弱小暗光目标难感知的问题，借鉴生物视觉对亮度变化、运动轮廓和多时间尺度刺激的敏锐响应机制，研究以事件相机为核心的极限场景感知方法。具体包括：研究事件触发的高速时空响应与目标检测定位，融合异构传感器在时间上的一致性和空间上的互补性，实现高速目标毫秒级检测、连续跟踪与三维定位；研究事件边缘空间驱动的精细表征方法，利用事件信号的时空边缘、维度扩展和短程双极性等特征，实现微小目标和细线障碍物的亚像素级检测；研究双时间尺度脉冲协同的极端环境感知增强方法，联合短时瞬态信息与长时稳定结构，抑制雨雪噪声并增强弱光目标响应。该部分从时间、空间和动态范围三个维度提升单点感知能力，使机器在极端场景中“看得更清”\footnote{公开支撑论文包括mmE-Loc、$x^2$-Fusion与PRE-Mamba等，详见“参考文献与资料来源”。}。

该研究内容内部形成由“及时发现”到“精细辨认”再到“稳定增强”的递进关系。高速响应研究首先保留传统帧式采样容易遗漏的瞬时运动过程，为目标跟踪和状态估计提供连续时间线索；精细表征研究进一步挖掘事件边缘与微弱运动之间的关联，使少像素、弱纹理目标仍能形成可分辨结构；双时间尺度研究则把短时变化与长时稳定信息结合起来，在保留敏捷响应的同时抑制偶发噪声与环境扰动。三类方法共同探索事件视觉从新型传感器走向可用感知系统所需的表示、算法与融合机制，为后续多节点协同提供更紧凑、更具时效性的前端信息。

\textbf{（2）多感官互补启发的全域协同融合。}
针对智慧城市基础设施中感知手段单一、通信带宽受限和采信判据缺失等问题，借鉴生物多感官系统的优势互补、选择传导和先验调节机制，研究多源、异构、分布式感知信息的协同融合方法。具体包括：研究多模态感官互补的异构融合方法，协同利用事件视觉的高时间分辨率、惯性信息的连续运动约束和三维空间结构的全局参考，构建跨模态时空匹配与层级融合机制；研究跨时空选择传导的多节点协同融合方法，根据感知质量、任务相关性和节点关系筛选高价值信息，并通过语义压缩和统计压缩降低传输冗余；研究先验知识调节增强的多源可信融合方法，引入可靠性以及语义、结构、运动等先验，对低可信、偏移和失真观测进行动态筛选、采信与校正。该部分沿“异构信息互补---高价值信息传导---多源信息可信调节”的路径提升全域观测的完备性、实时性和可靠性，使机器在大规模城市空间中“看得更广”\footnote{公开支撑论文包括EV-Pose、Communication-Efficient Collaborative Perception with Semantic and Statistical Compression与QUIDS，详见“参考文献与资料来源”。}。

该研究内容所解决的并不是一般意义上的数据拼接，而是多源信息从“能够对应”到“值得传输”再到“可以采信”的全过程。跨模态融合建立共同的时空参照，回答不同传感器观测是否描述同一对象、同一状态；选择传导根据任务需求压缩和分发信息，回答有限链路中应当优先传递什么；可信融合利用先验知识和动态可靠性判断信息质量，回答发生冲突时系统应当相信谁。三者相互依赖：没有对齐，信息价值无法准确比较；没有选择，规模化协同难以满足实时要求；没有可信调节，更多数据也可能放大误差。由此形成可随场景、节点和任务变化而调整的分布式协同机制。

\textbf{（3）认知卸载启发的多层次推理。}
针对交通调度、生产协同、安防预警和应急处置等任务中复杂语义理解、复杂空间认知与复杂系统规划难以兼顾准确性和实时性的问题，借鉴生物认知系统中快慢分流、分层记忆和高阶认知与灵敏执行协作的机制，研究大小模型、多层空间和异构智能体之间的任务卸载与协同推理方法。具体包括：研究快慢认知分流的复杂语义理解与按需推理，由轻量模型承担高频感知和状态判断任务，由基础大模型按需处理低频复杂语义和高阶决策任务；研究全局---局部空间分解的复杂空间认知与记忆规划，通过层级语义子目标、全局记忆与局部规划协同，降低长时程导航和探索的推理开销；研究高阶认知---灵敏执行协同的复杂系统规划，由大模型负责全局任务理解与协同规划，由异构智能体并行执行，并通过视觉词元筛选和弹性剪枝压缩冗余计算。该部分形成“任务分流---空间分解---规划执行协同”的多层次推理体系，使机器对复杂城市场景“看得更透”\footnote{公开支撑论文包括Embodied-R、CityNavAgent、MR-COGraphs与Balanced Token Pruning，详见“参考文献与资料来源”。}。

该研究内容以任务结构而不是模型大小作为资源分配依据。快慢认知分流先判断当前问题是否需要高阶语义推理，避免把所有输入无差别送入基础大模型；全局---局部空间分解把长时程任务转换为层级目标与可执行局部步骤，避免每次行动都重新处理完整环境；高阶认知与灵敏执行协同进一步把规划能力与平台执行能力匹配起来，使不同智能体围绕共享目标并行工作。分层记忆、视觉信息筛选和反馈更新贯穿其中，用于保持任务连续性、控制输入规模并根据执行结果修正规划，从而形成能够随任务复杂度动态调整的认知闭环。

\subsection{技术体系的内在联系与应用价值}

上述三项研究并非彼此孤立，而是围绕城市感知信息的生命周期形成逐层递进的技术体系：动态视觉机制首先提升极端条件下有效信息的获取能力，多感官互补机制进一步提升多源信息的覆盖范围、传导效率与融合可信度，认知卸载机制最终将融合后的信息转化为复杂任务中的语义理解、空间推理与协同决策。前一层为后一层提供更有效、更紧凑的信息输入，后一层的任务需求又可以反向调节前一层的信息筛选与资源分配，最终构成“感知---协同---认知---行动---反馈”的闭环。

这一闭环体现了本项目区别于单点算法改进的系统特征。在信息入口处，系统关注的不是尽可能采集全部数据，而是在物理约束下保留与任务相关的关键变化；在信息网络中，系统关注的不是无差别连接全部节点，而是建立跨模态、跨区域和跨主体的有效协作；在认知出口处，系统关注的不是把所有问题交给同一模型，而是使不同能力层级承担与其相匹配的任务。由此，感知、通信与计算资源可以围绕同一任务目标协调配置，局部技术进步也能够通过标准化的状态表达、价值评价和反馈接口传递到完整系统。

从应用角度看，极限场景感知可为高速交通参与者、低空飞行器、工业微小缺陷和暗光安防目标提供更及时的观测基础；全域协同融合可服务车路协同、园区巡检、多机搜索和跨区域应急等需要多节点共同覆盖的任务；多层次推理则可支撑自然语言任务理解、长时程导航、异构平台协作与动态规划。三类能力既可分别形成面向具体设备和场景的功能模块，也可组合为城市级智能感知平台，使研究成果具备从核心方法到系统集成、从实验验证到行业应用逐步转化的条件。

由此，本项目形成从单点极限感知、全域协同融合到多层次认知推理的完整技术链条，并面向智慧交通、智能制造、智能安防和智慧应急开展系统平台建设与产业应用。其预期贡献在于：以生物启发的信息处理原则突破传统均匀采样、全量传输和集中计算的局限，建立面向复杂开放环境的感知新方法；以跨层协同设计连接传感器、算法、网络、模型和智能体，提升技术在真实资源约束下的可用性；以典型城市任务牵引科研成果落地，为城市高质量发展、高效能治理和高水平安全提供关键技术支撑。

\clearpage
